% ------------------------------------------------------------------------------
% Formatvorlage für RFCCC1
% ------------------------------------------------------------------------------
%   erstellt von Christian Schiffler, 25.11.2010
%   basierend auf Arbeiten von Stefan Macke, 24.04.2009
%   http://blog.stefan-macke.de

% Dokumentenkopf ---------------------------------------------------------------
%   Diese Vorlage basiert auf "scrbook" aus dem koma-script.
% ------------------------------------------------------------------------------
\documentclass[
    fontsize=11pt, % Schriftgröße
    DIV10,
    a4paper, % Papierformat
    %oneside, % einseitiges Dokument
    openany, % remove empty pages after a chapter if chapter starts on the "left" hand
    titlepage, % es wird eine Titelseite verwendet
    parskip=half, % Abstand zwischen Absätzen (halbe Zeile)
    headings=normal, % Größe der Überschriften verkleinern
    listof=totoc, % Verzeichnisse im Inhaltsverzeichnis aufführen
    bibliography=totoc, % Literaturverzeichnis im Inhaltsverzeichnis aufführen
    index=totoc, % Index im Inhaltsverzeichnis aufführen
    captions=tableheading, % Beschriftung von Tabellen unterhalb ausgeben
    final % Status des Dokuments (final/draft)
]{scrbook}

% Umlaute ----------------------------------------------------------------------
%   Umlaute/Sonderzeichen wie äüöß direkt im Quelltext verwenden (CodePage).
%   Erlaubt automatische Trennung von Worten mit Umlauten.
% ------------------------------------------------------------------------------
\usepackage[utf8x]{inputenc}
\usepackage[T1]{fontenc}
\usepackage{textcomp} % Euro-Zeichen etc.

% Meta-Informationen -----------------------------------------------------------
%   Informationen über das Dokument, wie z.B. Titel, Autor, etc.
%   werden in der Datei Meta.tex definiert und können danach global
%   verwendet werden.
% ------------------------------------------------------------------------------
% Meta-Informationen -----------------------------------------------------------
%   Definition von globalen Parametern, die im gesamten Dokument verwendet
%   werden können (z.B auf dem Deckblatt etc.).
%
%   ACHTUNG: Wenn die Texte Umlaute oder ein Esszet enthalten, muss der folgende
%            Befehl bereits an dieser Stelle aktiviert werden:
%            \usepackage[latin1]{inputenc}
% ------------------------------------------------------------------------------
\newcommand{\titel}{Request for Contao Coder Comments 1 - RFCCC1}
\newcommand{\untertitel}{ Contao Extension Repository 3}
\newcommand{\autor}{Tristan Lins, Christian Schiffler}
\newcommand{\jahr}{2011}
\newcommand{\logo}{images/base/Title}
\newcommand{\deckblattlogo}{images/base/Title}


% benötigte Packages -----------------------------------------------------------
%   LaTeX-Packages, die benötigt werden, sind in die Datei Packages.tex
%   "ausgelagert", um diese Vorlage möglichst übersichtlich zu halten.
% ------------------------------------------------------------------------------
% Anpassung des Seitenlayouts --------------------------------------------------
%   siehe Seitenstil.tex
% ------------------------------------------------------------------------------
\usepackage[
    automark, % Kapitelangaben in Kopfzeile automatisch erstellen
    headsepline, % Trennlinie unter Kopfzeile
    ilines % Trennlinie linksbündig ausrichten
]{scrpage2}

\usepackage{ucs}

% Schrift ----------------------------------------------------------------------
\usepackage{lmodern} % bessere Fonts
\usepackage{relsize} % Schriftgröße relativ festlegen

% Befehle aus AMSTeX für mathematische Symbole z.B. \boldsymbol \mathbb --------
\usepackage{amsmath,amsfonts}

% für Index-Ausgabe mit \printindex --------------------------------------------
\usepackage{makeidx}

% für Glossar-Ausgabe mit \printglossaries -------------------------------------
\usepackage{glossaries}

% Einfache Definition der Zeilenabstände und Seitenränder etc. -----------------
\usepackage{setspace}
\usepackage{geometry}

% Symbolverzeichnis ------------------------------------------------------------
%   Symbolverzeichnisse bequem erstellen. Beruht auf MakeIndex:
%     makeindex.exe %Name%.nlo -s nomencl.ist -o %Name%.nls
%   erzeugt dann das Verzeichnis. Dieser Befehl kann z.B. im TeXnicCenter
%   als Postprozessor eingetragen werden, damit er nicht ständig manuell
%   ausgeführt werden muss.
%   Die Definitionen sind ausgegliedert in die Datei "Glossar.tex".
% ------------------------------------------------------------------------------
\usepackage[intoc]{nomencl}
\let\abbrev\nomenclature
\renewcommand{\nomname}{Abkürzungsverzeichnis}
\setlength{\nomlabelwidth}{.25\hsize}
\renewcommand{\nomlabel}[1]{#1 \dotfill}
\setlength{\nomitemsep}{-\parsep}

% zum Umfließen von Bildern ----------------------------------------------------
\usepackage{floatflt}


% zum Einbinden von Programmcode -----------------------------------------------
\usepackage{listings}
%\usepackage{scrhack}
\usepackage{xcolor} 
\definecolor{hellgelb}{rgb}{1,1,0.9}
\definecolor{colKeys}{rgb}{0,0,1}
\definecolor{colIdentifier}{rgb}{0,0,0}
\definecolor{colComments}{rgb}{1,0,0}
\definecolor{colString}{rgb}{0,0.5,0}
\lstset{
    float=hbp,
    basicstyle=\ttfamily\color{black}\small\smaller,
    identifierstyle=\color{colIdentifier},
    keywordstyle=\color{colKeys},
    stringstyle=\color{colString},
    commentstyle=\color{colComments},
    columns=flexible,
    tabsize=2,
    frame=single,
    extendedchars=true,
    showspaces=false,
    showstringspaces=false,
    numbers=left,
    numberstyle=\tiny,
    breaklines=true,
    backgroundcolor=\color{hellgelb},
    breakautoindent=true
}

% URL verlinken, lange URLs umbrechen etc. -------------------------------------
\usepackage{url}

% wichtig für korrekte Zitierweise ---------------------------------------------
\usepackage[square]{natbib}

% PDF-Optionen -----------------------------------------------------------------
\usepackage[
    bookmarks,
    bookmarksopen=true,
    colorlinks=true,
% diese Farbdefinitionen zeichnen Links im PDF farblich aus
    linkcolor=red, % einfache interne Verknüpfungen
    anchorcolor=black,% Ankertext
    citecolor=blue, % Verweise auf Literaturverzeichniseinträge im Text
    filecolor=magenta, % Verknüpfungen, die lokale Dateien öffnen
    menucolor=red, % Acrobat-Menüpunkte
    urlcolor=cyan, 
% diese Farbdefinitionen sollten für den Druck verwendet werden (alles schwarz)
    %linkcolor=black, % einfache interne Verknüpfungen
    %anchorcolor=black, % Ankertext
    %citecolor=black, % Verweise auf Literaturverzeichniseinträge im Text
    %filecolor=black, % Verknüpfungen, die lokale Dateien öffnen
    %menucolor=black, % Acrobat-Menüpunkte
    %urlcolor=black, 
    backref,
    plainpages=false, % zur korrekten Erstellung der Bookmarks
    pdfpagelabels, % zur korrekten Erstellung der Bookmarks
    hypertexnames=false, % zur korrekten Erstellung der Bookmarks
    linktocpage % Seitenzahlen anstatt Text im Inhaltsverzeichnis verlinken
]{hyperref}
% Befehle, die Umlaute ausgeben, führen zu Fehlern, wenn sie hyperref als Optionen übergeben werden
\hypersetup{
    pdftitle={\titel \untertitel},
    pdfauthor={\autor},
    pdfcreator={\autor},
    pdfsubject={\titel \untertitel},
    pdfkeywords={\titel \untertitel},
}

% Grafiken ---------------------------------------------------------------------

% for UML
\usepackage{emp}
% are we pdflatex?
\ifx\pdftexversion\undefined
% allow usage of jpeg, mps etc. in non pdf mode
\usepackage[dvips,final]{graphicx}
%enable landscape
\usepackage{lscape}
\else
% allow usage of jpeg, mps etc. in pdf mode
\usepackage[dvips,final,pdftex]{graphicx}
\DeclareGraphicsRule{*}{mps}{*}{}
%enable landscape
\usepackage{pdflscape}
\fi

% hier liegen die Bilder des Dokuments
\graphicspath{{images/}}


% fortlaufendes Durchnummerieren der Fußnoten ----------------------------------
\usepackage{chngcntr}

% für lange Tabellen -----------------------------------------------------------
\usepackage{longtable}
\usepackage{array}
\usepackage{ragged2e}
\usepackage{lscape}

% Spaltendefinition rechtsbündig mit definierter Breite ------------------------
\newcolumntype{w}[1]{>{\raggedleft\hspace{0pt}}p{#1}}

% Formatierung von Listen ändern -----------------------------------------------
\usepackage{paralist}

% bei der Definition eigener Befehle benötigt
\usepackage{ifthen}

% definiert u.a. die Befehle \todo und \listoftodos
\usepackage{todonotes}

% sorgt dafür, dass Leerzeichen hinter parameterlosen Makros nicht als Makroendezeichen interpretiert werden
\usepackage{xspace}

% For floating of images within text
\usepackage{wrapfig}

% for underline via \uline, \uwave and \sout
\usepackage[normalem]{ulem}

% http://mirror.ctan.org/macros/generic/dirtree.zip
\usepackage{dirtree}

% Erstellung eines Index und Abkürzungsverzeichnisses aktivieren ---------------
\makeindex
\makenomenclature
\makeglossaries

% Kopf- und Fußzeilen, Seitenränder etc. ---------------------------------------
% Zeilenabstand 1,5 Zeilen -----------------------------------------------------
%\onehalfspacing
\spacing{1.0}

% Seitenränder -----------------------------------------------------------------
\setlength{\topskip}{\ht\strutbox} % behebt Warnung von geometry
\geometry{paper=a4paper}

% Kopf- und Fußzeilen ----------------------------------------------------------
\pagestyle{scrheadings}
% Kopf- und Fußzeile auch auf Kapitelanfangsseiten
\renewcommand*{\chapterpagestyle}{scrheadings} 
% Schriftform der Kopfzeile
\renewcommand{\headfont}{\normalfont}

% Kopfzeile
\ihead{\large{\textsc{\titel}}\\ \small{\untertitel} \\[2ex] \textit{\headmark}}
\chead{}
%\ohead{\includegraphics[scale=0.075]{\logo}}
\setlength{\headheight}{21mm} % Höhe der Kopfzeile
% Kopfzeile über den Text hinaus verbreitern
%\setheadwidth[0pt]{textwithmarginpar} 
\setheadsepline[text]{0.4pt} % Trennlinie unter Kopfzeile

% Fußzeile
\ifoot{\copyright\ \autor}
\cfoot{}
\ofoot{\pagemark}

% sonstige typographische Einstellungen ----------------------------------------

% erzeugt ein wenig mehr Platz hinter einem Punkt
\frenchspacing 

% Schusterjungen und Hurenkinder vermeiden
\clubpenalty = 10000
\widowpenalty = 10000 
\displaywidowpenalty = 10000

% Quellcode-Ausgabe formatieren
\lstset{numbers=left, numberstyle=\tiny, numbersep=5pt, breaklines=true}
\lstset{emph={square}, emphstyle=\color{red}, emph={[2]root,base}, emphstyle={[2]\color{blue}}}

% Fußnoten fortlaufend durchnummerieren
\counterwithout{footnote}{chapter}

%% ... und nimm alle 4 Ebenen in das Inhaltsverzeichnis auf.
\setcounter{tocdepth}{4}

% eigene Definitionen für Silbentrennung
% Trennvorschläge im Text werden mit \" angegeben
% untrennbare Wörter und Ausnahmen von der normalen Trennung können in dieser
% Datei mittels \hyphenation definiert werden


% eigene LaTeX-Befehle
% Eigene Befehle und typographische Auszeichnungen für diese

% einfaches Wechseln der Schrift, z.B.: \changefont{cmss}{sbc}{n}
\newcommand{\changefont}[3]{\fontfamily{#1} \fontseries{#2} \fontshape{#3} \selectfont}

% Abkürzungen mit korrektem Leerraum 
\newcommand{\ua}{\mbox{u.\,a.\ }}
\newcommand{\zB}{\mbox{z.\,B.\ }}
\newcommand{\dahe}{\mbox{d.\,h.\ }}
\newcommand{\Vgl}{Vgl.\ }
\newcommand{\bzw}{bzw.\ }
\newcommand{\evtl}{evtl.\ }
\newcommand{\ie}{\mbox{i.\,e.\ }}

\newcommand{\illustration}[1]{illustration~\ref{fig:#1}}

\newcommand{\bs}{$\backslash$}

\newcommand{\sectiontitle}{}
\let\TeXsection\section
\renewcommand{\section}[2][\empty]{\ifthenelse{\equal{#1}{\empty}}{\TeXsection{#2}\renewcommand{\sectiontitle}{#2}}{\TeXsection{#2}\label{#1}\renewcommand{\sectiontitle}{#1}}}

\newcommand{\subsectiontitle}{}
\let\TeXsubsection\subsection
\renewcommand{\subsection}[2][\empty]{\ifthenelse{\equal{#1}{\empty}}{\TeXsubsection{#2}\renewcommand{\subsectiontitle}{#2}}{\TeXsubsection{#2}\label{#1}\renewcommand{\subsectiontitle}{#1}}}

\let\TeXsubsubsection\subsubsection
\renewcommand{\subsubsection}[2][\empty]{\ifthenelse{\equal{#1}{\empty}}{\TeXsubsubsection{#2}}{\TeXsubsubsection{#2}\label{#1}}}

\newcommand{\supref}[1]{\textsuperscript{\ref{#1}}}

% erzeugt ein Listenelement mit fetter Überschrift 
\newcommand{\itemd}[2]{\item{\textbf{#1}}\\{#2}}

% einige Befehle zum Zitieren --------------------------------------------------
\newcommand{\Zitat}[2][\empty]{\ifthenelse{\equal{#1}{\empty}}{\citep{#2}}{\citep[#1]{#2}}}

% zum Ausgeben von Autoren
\newcommand{\AutorName}[1]{\textsc{#1}}
\newcommand{\Autor}[1]{\AutorName{\citeauthor{#1}}}

% verschiedene Befehle um Wörter semantisch auszuzeichnen ----------------------
\newcommand{\newDef}[1]{\textbf{#1}}
\newcommand{\techTerm}[1]{\textit{#1}}

\newcommand{\DisplInput}[1]{\texttt{#1}}
\newcommand{\DisplCode}[1]{\texttt{#1}}
\newcommand{\DisplFile}[1]{\texttt{#1}}

\newcommand{\Datentyp}[1]{\textsf{#1}}
\newcommand{\XMLElement}[1]{\textsf{#1}}
\newcommand{\Webservice}[1]{\textsf{#1}}

% Deklaration von Hinweistexten
\newcommand{\hint}[1]{
\begin{tabular}{w{1.5cm}p{12cm}}\\
\scalebox{0.125}{\includegraphics{images/base/Hint_Notification}} & #1
\end{tabular}
}

% Deklaration von Warnungen
\newcommand{\warning}[1]{
\begin{tabular}{w{1.5cm}p{12cm}}\\
\scalebox{0.125}{\includegraphics{images/base/Warning_Notification}} & #1
\end{tabular}
}

% einbinden eines emp Bildes.
\newcommand{\useempcenter}[2]{\begin{figure}[h!]\centering\scalebox{0.8}{\empuse{#1}}\caption{#2}\end{figure}}

\newcommand{\useempfloat}[3]{\begin{wrapfigure}{#1}{0.4\textwidth}\begin{center}\empuse{#2}\end{center}\caption{#3}\end{wrapfigure}}

\newcommand{\useemp}[3][\empty]{\ifthenelse{\equal{#1}{\empty}}{\useempcenter{#2}{#3}}{\useempfloat{#1}{#2}{#3}}}

% Wrapper für useemp für uml Klassendiagramme
\newcommand{\umlclass}[3][\empty]{\ifthenelse{\equal{#1}{\empty}}{\useemp{#2}{#3}}{\useemp[#1]{#2}{#3}}}

% Wrapper um listings einzubinden
\newcommand{\phplisting}[3]{\lstset{language=PHP, basicstyle=\footnotesize, showstringspaces=false, tabsize=2}\lstinputlisting[label=#1,caption=#2]{#3}}
\newcommand{\xmllisting}[3]{\lstset{language=XML, basicstyle=\footnotesize, showstringspaces=false, tabsize=2}\lstinputlisting[label=#1,caption=#2]{#3}}


% Das eigentliche Dokument -----------------------------------------------------
%   Der eigentliche Inhalt des Dokuments beginnt hier. Die einzelnen Seiten
%   und Kapitel werden in eigene Dateien ausgelagert und hier nur inkludiert.
% ------------------------------------------------------------------------------
\begin{document}

% auch subsubsection nummerieren
\setcounter{secnumdepth}{3}
\setcounter{tocdepth}{3}

% Deckblatt und Abstract ohne Seitenzahl
\ofoot{}
\thispagestyle{plain}
\begin{titlepage}

\begin{center}

\huge{\textbf{\titel}}\\[1.5ex]
\LARGE{\textbf{\untertitel}}\\[6ex]
%\includegraphics[scale=0.2]{\deckblattlogo}\\[6ex]

\normalsize
\autor\\[1.2ex]
\copyright\ \jahr\\[9ex]

\end{center}

\singlespacing
\small
\noindent 
\begin{center}
This work is licensed under a \\ \hyperlink{http://creativecommons.org/licenses/by-sa/3.0/}{Creative Commons Attribution-ShareAlike 3.0 Unported License}.
\end{center}

A human-readable summary of the Legal Code (the full license):

You are free:
\begin{itemize}
\item to Share — to copy, distribute and transmit the work
\item to Remix — to adapt the work
\item to make commercial use of the work
\end{itemize}

Under the following conditions:
\begin{itemize}
\item Attribution — You must attribute the work in the manner specified by the author or licensor (but not in any way that suggests that they endorse you or your use of the work).
\item Share Alike — If you alter, transform, or build upon this work, you may distribute the resulting work only under the same or similar license to this one.
\end{itemize}

Read the license for further details.

\end{titlepage}

\chapter{Abstract}

This document serves as a discussion and documentation foundation for the new Extension Repository version 3 which shall replace the existing repository for Contao.

This document will be proposed for discussion at Contao Camp 2011.

The first goal of this document is, to identify problems with the current extension repository and provide solutions to them.

To fulfil this goal, the following conditions have to be met:
\begin{itemize}
\item Compatibility to Contao version 2 and version 3 (upcoming).
\item Be more flexible than the current system.
\item (almost) anyone shall be able to run his/her own repository.
\item Installation and uninstallation routines to aid extensions in setup.
\item Standardized package format to ease installation.
\item Improved dependencies.
\item Add possibility to define conflicting extensions.
\end{itemize}

The open discussion on this document is currently taking place in the 
\href{http://www.contao-community.de/showthread.php?22031-Gruppe-neuer-ER-Client-und-ER3}{German Contao Community}

\warning{This document is still a work in progress and should not be used for implementation yet as all parts are still subject to change. Once the document is finalized, it will get a version number and will be released along with some example implementation of the services.}
\ofoot{\pagemark}


% arabische Seitenzahlen im Hauptteil ------------------------------------------
\clearpage
\pagenumbering{arabic}
\tableofcontents % Inhaltsverzeichnis

% UML support
\begin{empfile}
\begin{empcmds}
input metauml;
\end{empcmds}
% die Inhaltskapitel werden in "Content.tex" inkludiert -------------------------
% Hier können die einzelnen Kapitel inkludiert werden. Sie müssen in den 
% entsprechenden .TEX-Dateien vorliegen. Die Dateinamen können natürlich 
% angepasst werden.

\chapter{Changelog}

\newcommand{\ChangeLogEntry}[3]{\subparagraph{Changes by #1 on #2} \begin{itemize} #3 \end{itemize} }

\tiny
\ChangeLogEntry{Tristan Lins}{08.09.2011}{
	\item Initial draft containing document description and package definition.
}
\ChangeLogEntry{Tristan Lins}{09.09.2011}{
	\item added PACKAGE XML information file.
	\item small adjustments in PACKAGE INI Format.
	\item elaborated usage of versions with an extension component.
}
\ChangeLogEntry{Tristan Lins}{10.09.2011}{
	\item Added first description of ER3 Repository, ER3 Repository Manager and ER3 Portal.
}
\ChangeLogEntry{Tristan Lins}{11.09.2011}{
	\item Removed INI Meta file specification.
	\item Reworked PACKAGE information file.
	\item Reworked XML schema.
	\item Reworked XML example.
}
\ChangeLogEntry{Tristan Lins}{12.09.2011}{
	\item Reworked XML Schema and example.
	\item added "versioning of dependencies" and "up/downgrade compatibility" to PACKAGE information.
	\item Removed description of license file and EULA.
}
\ChangeLogEntry{Christian Schiffler}{13.09.2011-14.09.2011}{
	\item Transition to \LaTeX{}, 
	\item translation to English and major rephrasing of all sections
	\item removed \techTerm{lowest\_bound} and \techTerm{highest\_bound} from \DisplFile{extension.xsd}.
}
\ChangeLogEntry{Tristan Lins}{15.09.2011}{
	\item Reworked Repository Portal\supref{sec:repository portal} section
}
\ChangeLogEntry{Christian Schiffler}{16.09.2011}{
	\item finished translation to English (mainly repository\supref{sec:repository} and repository manager \supref{sec:repository manager}).
}
\ChangeLogEntry{Tristan Lins}{19.09.2011}{
	\item Removed package scrhack, it generates the firstpage with "stylistings".
	\item Added first description of "Installation"\supref{sec:installation}, "Installation Process"\supref{sec:installation process} and "Installer"\supref{sec:installer_interfaces_and_classes}.
}
\ChangeLogEntry{Tristan Lins}{25.09.2011}{
	\item Rework complete chapter "Installation"
	\item Added section "ER3 Manager"\supref{sec:er3 manager}.
	\item Added section "Contao Installation Manager"\supref{sec:conatao installation manager}.
	\item Rename section "Installer" to "Installer interfaces and classes"\supref{sec:installer interfaces and classes}.
	\item Added subsection "InstallerInformation"\supref{sec:installerinformation}.
}
\ChangeLogEntry{Tristan Lins}{03.10.2011}{
	\item Added subsection "Extension Builder"\supref{sec:extension builder} to section "Repository Portal"\supref{sec:repository portal}
	\item Added section "Namespaces"\supref{sec:namespaces}.
}
\ChangeLogEntry{Tristan Lins}{12.10.2011}{
	\item Added subsection "Comments and rating"\supref{sec:comments and rating} to section "Repository Portal"\supref{sec:repository portal}
}
\ChangeLogEntry{Tristan Lins}{16.10.2011}{
	\item Added encrypted flag to "PACKAGE information file schema"\supref{sec:package information file schema}
	\item Added chapter Licensing
	\item Added appendix "Unencrypted package checksum calculation (upc.php)"\supref{sec:unencrypted package checksum calculation}
	\item Added appendix "Package encryption (encrypt.php)"\supref{sec:package encryption implementation}
	\item Added appendix "Installation id calculation (installationid.php)"\supref{sec:unencrypted package checksum calculation}
	\item Added example "Hello World!" extension package
}

\normalsize
% Anpassung des Seitenlayouts --------------------------------------------------
%   siehe Seitenstil.tex
% ------------------------------------------------------------------------------
\usepackage[
    automark, % Kapitelangaben in Kopfzeile automatisch erstellen
    headsepline, % Trennlinie unter Kopfzeile
    ilines % Trennlinie linksbündig ausrichten
]{scrpage2}

\usepackage{ucs}

% Schrift ----------------------------------------------------------------------
\usepackage{lmodern} % bessere Fonts
\usepackage{relsize} % Schriftgröße relativ festlegen

% Befehle aus AMSTeX für mathematische Symbole z.B. \boldsymbol \mathbb --------
\usepackage{amsmath,amsfonts}

% für Index-Ausgabe mit \printindex --------------------------------------------
\usepackage{makeidx}

% für Glossar-Ausgabe mit \printglossaries -------------------------------------
\usepackage{glossaries}

% Einfache Definition der Zeilenabstände und Seitenränder etc. -----------------
\usepackage{setspace}
\usepackage{geometry}

% Symbolverzeichnis ------------------------------------------------------------
%   Symbolverzeichnisse bequem erstellen. Beruht auf MakeIndex:
%     makeindex.exe %Name%.nlo -s nomencl.ist -o %Name%.nls
%   erzeugt dann das Verzeichnis. Dieser Befehl kann z.B. im TeXnicCenter
%   als Postprozessor eingetragen werden, damit er nicht ständig manuell
%   ausgeführt werden muss.
%   Die Definitionen sind ausgegliedert in die Datei "Glossar.tex".
% ------------------------------------------------------------------------------
\usepackage[intoc]{nomencl}
\let\abbrev\nomenclature
\renewcommand{\nomname}{Abkürzungsverzeichnis}
\setlength{\nomlabelwidth}{.25\hsize}
\renewcommand{\nomlabel}[1]{#1 \dotfill}
\setlength{\nomitemsep}{-\parsep}

% zum Umfließen von Bildern ----------------------------------------------------
\usepackage{floatflt}


% zum Einbinden von Programmcode -----------------------------------------------
\usepackage{listings}
%\usepackage{scrhack}
\usepackage{xcolor} 
\definecolor{hellgelb}{rgb}{1,1,0.9}
\definecolor{colKeys}{rgb}{0,0,1}
\definecolor{colIdentifier}{rgb}{0,0,0}
\definecolor{colComments}{rgb}{1,0,0}
\definecolor{colString}{rgb}{0,0.5,0}
\lstset{
    float=hbp,
    basicstyle=\ttfamily\color{black}\small\smaller,
    identifierstyle=\color{colIdentifier},
    keywordstyle=\color{colKeys},
    stringstyle=\color{colString},
    commentstyle=\color{colComments},
    columns=flexible,
    tabsize=2,
    frame=single,
    extendedchars=true,
    showspaces=false,
    showstringspaces=false,
    numbers=left,
    numberstyle=\tiny,
    breaklines=true,
    backgroundcolor=\color{hellgelb},
    breakautoindent=true
}

% URL verlinken, lange URLs umbrechen etc. -------------------------------------
\usepackage{url}

% wichtig für korrekte Zitierweise ---------------------------------------------
\usepackage[square]{natbib}

% PDF-Optionen -----------------------------------------------------------------
\usepackage[
    bookmarks,
    bookmarksopen=true,
    colorlinks=true,
% diese Farbdefinitionen zeichnen Links im PDF farblich aus
    linkcolor=red, % einfache interne Verknüpfungen
    anchorcolor=black,% Ankertext
    citecolor=blue, % Verweise auf Literaturverzeichniseinträge im Text
    filecolor=magenta, % Verknüpfungen, die lokale Dateien öffnen
    menucolor=red, % Acrobat-Menüpunkte
    urlcolor=cyan, 
% diese Farbdefinitionen sollten für den Druck verwendet werden (alles schwarz)
    %linkcolor=black, % einfache interne Verknüpfungen
    %anchorcolor=black, % Ankertext
    %citecolor=black, % Verweise auf Literaturverzeichniseinträge im Text
    %filecolor=black, % Verknüpfungen, die lokale Dateien öffnen
    %menucolor=black, % Acrobat-Menüpunkte
    %urlcolor=black, 
    backref,
    plainpages=false, % zur korrekten Erstellung der Bookmarks
    pdfpagelabels, % zur korrekten Erstellung der Bookmarks
    hypertexnames=false, % zur korrekten Erstellung der Bookmarks
    linktocpage % Seitenzahlen anstatt Text im Inhaltsverzeichnis verlinken
]{hyperref}
% Befehle, die Umlaute ausgeben, führen zu Fehlern, wenn sie hyperref als Optionen übergeben werden
\hypersetup{
    pdftitle={\titel \untertitel},
    pdfauthor={\autor},
    pdfcreator={\autor},
    pdfsubject={\titel \untertitel},
    pdfkeywords={\titel \untertitel},
}

% Grafiken ---------------------------------------------------------------------

% for UML
\usepackage{emp}
% are we pdflatex?
\ifx\pdftexversion\undefined
% allow usage of jpeg, mps etc. in non pdf mode
\usepackage[dvips,final]{graphicx}
%enable landscape
\usepackage{lscape}
\else
% allow usage of jpeg, mps etc. in pdf mode
\usepackage[dvips,final,pdftex]{graphicx}
\DeclareGraphicsRule{*}{mps}{*}{}
%enable landscape
\usepackage{pdflscape}
\fi

% hier liegen die Bilder des Dokuments
\graphicspath{{images/}}


% fortlaufendes Durchnummerieren der Fußnoten ----------------------------------
\usepackage{chngcntr}

% für lange Tabellen -----------------------------------------------------------
\usepackage{longtable}
\usepackage{array}
\usepackage{ragged2e}
\usepackage{lscape}

% Spaltendefinition rechtsbündig mit definierter Breite ------------------------
\newcolumntype{w}[1]{>{\raggedleft\hspace{0pt}}p{#1}}

% Formatierung von Listen ändern -----------------------------------------------
\usepackage{paralist}

% bei der Definition eigener Befehle benötigt
\usepackage{ifthen}

% definiert u.a. die Befehle \todo und \listoftodos
\usepackage{todonotes}

% sorgt dafür, dass Leerzeichen hinter parameterlosen Makros nicht als Makroendezeichen interpretiert werden
\usepackage{xspace}

% For floating of images within text
\usepackage{wrapfig}

% for underline via \uline, \uwave and \sout
\usepackage[normalem]{ulem}

% http://mirror.ctan.org/macros/generic/dirtree.zip
\usepackage{dirtree}

\chapter{Repository Stack}

In this chapter you find the whole extension repository stack defined.

\begin{landscape}
% I know, this is not really a state diagram but rather entities but hey.
\begin{figure}
\begin{emp}[Repository Stack](20, 20)

StateInfo.iRepositoryDistributionNetwork;
iRepositoryDistributionNetwork.foreColor := (.47, .94, .27);

NoteInfo.iUMLNote;
iUMLNote.foreColor := (0, .9, .9);


State.Portal("Repository Portal")();
ENote.NotePortal(iUMLNote)(
	"Public browsing interface",
	"- Public registration",
	"- Public profile management",
	"- Social networking");

State.RepManager("Repository Manager")();
ENote.NoteRepManager(iUMLNote)(
	"Authorization required management interface",
	"- User/Account management",
	"- Package builder",
	"- Package management",
	"- Administration");

State.RepMaster("Repository (Master)")();
ENote.NoteRepMaster(iUMLNote)("Package distribution network");

State.RepSlaveA("Repository (Slave1)")();
State.RepSlaveB("Repository (Slave2)")();
State.RepSlaveC("Repository (Slave3)")();

% this is ugly as hell but I was unable to position the notes more properly using
% same.left(NotePortal, NoteRepManager, NoteRepMaster);
% it always threw an error, FIX this.
leftToRight(95)(Portal, NotePortal);
leftToRight(90)(RepManager, NoteRepManager);
leftToRight(150)(RepMaster, NoteRepMaster);

% combine distribution network
leftToRight(20)(RepSlaveA, RepSlaveB);
topToBottom(20)(RepMaster, RepSlaveA);
EState.RepositoryDistributionNetwork(iRepositoryDistributionNetwork)("Package distribution network")(RepMaster, RepSlaveA, RepSlaveB);

% positioning of globals
topToBottom(20)(Portal, RepManager, RepositoryDistributionNetwork);

State.RepositoryStack("Repository Stack")(Portal, RepManager, RepositoryDistributionNetwork);

drawObjects(RepositoryStack, NotePortal, NoteRepManager, NoteRepMaster);

clink(transition)(RepMaster, RepSlaveA);
clink(transition)(RepMaster, RepSlaveB);

clink(transition)(Portal, RepManager);
clink(transition)(RepManager, RepMaster);


clink(dashedLink)(Portal, NotePortal);
clink(dashedLink)(RepManager, NoteRepManager);
clink(dashedLink)(RepMaster, NoteRepMaster);

\end{emp}
\caption{Repository Stack}
\end{figure}

\end{landscape}

\chapter{Repository}
\todo{translation}

Das Repository ist der Paketspeicher, bietet Zugriff auf die Paketinformationen und verteilt die Pakete. Für den Zugriff und die Verwaltung der Pakete im Repository steht eine REST Schnittstelle zur Verfügung. Das Repository verfügt demnach nicht über ein eigenes Web-Interface.
Paketspeicher

Im Paketspeicher werden die Pakete abgelegt. Der Paketspeicher ist strukturiert aufgebaut, daraus ergeben sich folgende Paketpfade:
/<erster Buchstabe>/<name>/<name>-<version>.cext
z.B.
/a/Avisota/Avisota-1.5.1.cext
Die Pfade werden direkt auf die Installationswurzel gemapt.
REST Schnittstelle

Die REST Schnittstelle dient zum Abrufen von Informationen (ohne Authentifizierung) und zum Verwalten des Repositories (mit Authentifizierung).
Sie kann über dem Pfad /rest/ abgefragt werden.
Als Suffix wird das Abfrageformat angehängt, z.B. /json
Vor dem Abfrageformat werden die Parameter als Pfad-Komponente eingefügt, z.B. /rest/.../parameter/value/json
Schnittstellen

/rest/list/recent
Auflisten der neuesten Pakete.
Erlaubte Parameter: filter, offset, limit
/rest/list/categories
Auflisten der Kategorien in dem Repository.
Erlaubte Parameter: offset, limit
/rest/list/tags
Auflisten der Tags in dem Repository.
Erlaubte Parameter: offset, limit
/rest/search/keywords
Das Repository nach dem Schlüsselwörtern keywords durchsuchen.
Erlaubte Parameter: filter, offset, limit
/rest/put
Ein Paket in das Repository laden. Das Paket wird als POST Daten übertragen.
/rest/push
Ein paket in das Repository laden. Im Gegensatz zur put Methode wird in den POST Daten eine URL übertragen, vor der aus das Repository das Paket beziehen kann.
/rest/delete/name[/version]
Ein Paket aus dem Repository löschen. Das Paket wird über den Paketnamen identifiziert. Wenn die version angegeben wird, dann wird nur die angegebene Version gelöscht, ansonsten werden alle Versionen gelöscht.
/rest/settings/name/version
Liefert die Einstellungen (Kategorie, Veröffentlichung, Zugriff) einer Erweiterung im Repository.
/rest/configure/name/version
Die Einstellungen einer Erweiterung neu setzen. Als PUSH Daten werden die neuen Einstellungen übergeben.
Parameter

filter

Nach Kategorien filtern. Komma getrennte Liste.
z.B. /rest/list/rectent/filter/plugin/json
offset

Die ersten Treffer der Ergebnisliste überspringen.
z.B. /rest/list/recent/offset/10/json
limit

Das Ergebnis auf eine gewisse Anzahl Treffer begrenzen.
z.B. /rest/list/recent/limit/10/json
Ergebnisobjekt

Authentifizierung

Zum Zugriff auf die REST Schnittstelle gibt es 2 Authentifizierungsmethoden.
Authentifizierung mit Benutzername+Passwort
Authentifizierung mit AuthKey
Alle Zugriffe lassen sich zusätzlich per IP Filtern.
Setup Interface

Das Setup Interface ist das einzige Web-Interface, es dient zur Installation, Aktualisierung, Wartung und Verwaltung der Zugriffsrechte.


\chapter{Repository Manager}
\todo{translation}
Der Repository Manager ist ein Web-Interface und die einfachste Form, ein Repository zu verwalten. Es bietet alle Basisfunktionen für das Repository. Es ist Teil des Portals, kann aber auch unabhängig davon genutzt werden.
Aus der Verbindung von Repository und Repository Manager lässt sich ein einfaches, leicht zu verwaltendes Repository aufbauen, das nur über die JSON Schnittstelle von außen durchsuchbar ist.
Diese Installation ist vor allem für Agenturen und private Repositories gedacht, die kein öffentliches Portal benötigen.


\section[sec:repository portal]{Repository Portal}

The portal is a public interface. It links external services to provide additional informations for extensions, that are available on the world wide web.

First it is a repository browser. Users can browse and search extensions and show details of them.
Secondary it is a information aggregator to display related informations.

A package contains links to other sources, see \ref{sec:package information file}.
\lstset{language=XML} 
\begin{lstlisting}[caption=XML example: links to other sources]
...
<ctx:link rel="homepage" href="http://www.contao.org">
    <ctx:description>Official project homepage.</ctx:description>
</ctx:link>
<ctx:link rel="manual" href="http://www.contao.org/documentation.html" language="en">
    <ctx:description>Official User-Manual.</ctx:description>
</ctx:link>
<ctx:link rel="help" href="http://en.contaowiki.org" language="en">
    <ctx:description>Inofficial Wiki.</ctx:description>
</ctx:link>
<ctx:link rel="video" href="http://www.youtube.com/watch?feature=player_embedded&amp;v=8HkQVmfjjYc">
    <ctx:description>Introduction into TypoLIGHT/Contao.</ctx:description>
</ctx:link>
<ctx:link rel="forum" href="http://www.contao-community.de" language="de">
    <ctx:description>German contao community.</ctx:description>
</ctx:link>
<ctx:link rel="tracker" href="http://dev.contao.org">
    <ctx:description>Official bug and issue tracker.</ctx:description>
</ctx:link>
<ctx:link rel="scm" href="http://svn.contao.org">
    <ctx:description>Official svn repository.</ctx:description>
</ctx:link>
...
\end{lstlisting}
The portal use this links and try to fetch related informations from these systems. I.e. this can be bugs and issues from a tracking system like Redmine or Github. Or screencast from YouTube. Or tutorials from a MediaWiki.

The Portal is not only an information aggregator. Not only the developer can link external informations to a extensions.
Other users can link ther own informations, like tutorials, screencasts or similar to an extension.
Theses links are stored in the portal, but not in the extension information. This is the social networking character of the Portal.

\subsection[sec:information kinds and system connectors]{Information kinds and system connectors}

The Portal use multiple connectors to get informations from other systems.
First we have to define the kind of informations, we are able to aggregate.

\begin{itemize}
\item Bugs and issues\\
The portal can display bugs and issues to show the current development status and upstream development.
\item Manuals and Tutorials - Wiki informations\\
One of the best ways to manage manuals and tutorials is to use a wiki. The portal can display related wiki articles.
\item Videos and Screencasts\\
Videos and screencasts are another kind of tutorials. The portal can display videos hosted on video streaming platforms like YouTube.
\item Source code\\
If an scm repository is configured, the portal can directly link to the source code.
\end{itemize}

These are just a few possible information kinds, the repository may use.

Next we have to define how the informations are aggregated.
A lot of systems like \hyperlink{http://www.mediawiki.org/wiki/API:Main_page/de}{MediaWiki}, \hyperlink{http://code.google.com/intl/de-DE/apis/youtube/overview.html}{YouTube} and \hyperlink{http://developer.github.com/v3/}{github} provide external APIs.
The Portal use connectors to this systems to get informations.
Because not all theses services are uniq, like MediaWiki the systems have to be configures dynamicaly.
All users are allowed to request adding a service, but only administrators are able to add them finaly into the portal.
Initialy there may be these connectors available:

\small
\begin{longtable}{|p{0.25\textwidth}|c|c|c|c}
\caption{Information system connectors and information kinds} \\
\hline
\label{tab:information system connectors and information kinds}
\textbf{System} & \textbf{Tracking} & \textbf{Wiki} & \textbf{Video} & \textbf{SCM}
\\ \hline
 MediaWiki & & X & &
\\ \hline
 Redmine & X & X & &
\\ \hline
 Github & X & X & & X
\\ \hline
 YouTube & & & X &
\\ \hline
 Vimeo & & & X &
\\ \hline
\end{longtable}
\normalsize


\chapter[sec:installation]{Installation}


\section[sec:installation process]{Installation Process}

This section describe the flow of the installation process.

\todo{Make some charts to visualize the processes}

\subsection[sec:install process]{Install Process}

\begin{enumerate}
\item Download the package into local package archive
\item Read the installer class from package meta
\item Instanciate the installer
\item Create new installer information object
\item Call the installers \DisplFile{install(..)} method
\item Store the installer information object
\end{enumerate}

\subsection[sec:upgrade process]{Upgrade Process}

\begin{enumerate}
\item Download the new version package into local package archive
\item Read the installer class from package meta
\item Instanciate the installer
\item Load installer information object
\item Call the installers \DisplFile{upgrade(..)} method
\item Store the installer information object
\end{enumerate}

\subsection[sec:downgrade process]{Downgrade Process}

\paragraph{Method 1: Upgrade is possible}

\begin{enumerate}
\item Download the old version package into local package archive
\item Read the installer class from current installed package meta
\item Instanciate the installer of the current installed version
\item Load installer information object
\item Call the installers \DisplFile{downgrade(..)} method
\item Read the installer class from old version package meta
\item Instanciate the installer of the old version
\item Call the installers \DisplFile{upgrade(..)} method
\item Store the installer information object
\end{enumerate}

\paragraph{Method 2: Upgrade not possible - safe method}

\begin{enumerate}
\item Download the old version package into local package archive
\item Read the installer class from current installed package meta
\item Instanciate the installer of the current installed version
\item Load installer information object
\item Call the installers \DisplFile{uninstall(..)} method
\item Read the installer class from old version package meta
\item Instanciate the installer of the old version
\item Call the installers \DisplFile{install(..)} method
\item Store the installer information object
\end{enumerate}

\paragraph{Method 3: Upgrade not possible - hard method}

\begin{enumerate}
\item Download the old version package into local package archive
\item Read the installer class from current installed package meta
\item Instanciate the installer of the current installed version
\item Load installer information object
\item Call the installers \DisplFile{uninstall(..)} method
\item Call the installers \DisplFile{clean(..)} method
\item Read the installer class from old version package meta
\item Instanciate the installer of the old version
\item Call the installers \DisplFile{install(..)} method
\item Store the installer information object
\end{enumerate}

\subsection[sec:reinstall process]{Reinstall Process}

\begin{enumerate}
\item Download the package into local package archive
\item Read the installer class from package meta
\item Instanciate the installer
\item Load installer information object
\item Call the installers \DisplFile{reinstall(..)} method
\item Store the installer information object
\end{enumerate}

\subsection[sec:uninstall process]{Uninstall Process}

\begin{enumerate}
\item Download the package into local package archive
\item Read the installer class from package meta
\item Instanciate the installer
\item Load installer information object
\item Call the installers \DisplFile{uninstall(..)} method
\item Store the installer information object (maybe some DATA files and schema is just remeaning in the system)
\end{enumerate}

\subsection[sec:clean uninstall process]{Clean Uninstall Process}

Same as clean process.

\subsection[sec:clean process]{Clean Process}

\begin{enumerate}
\item Download the package into local package archive
\item Read the installer class from package meta
\item Instanciate the installer
\item Load installer information object
\item If still installed, call the installers \DisplFile{uninstall(..)} method
\item Call the installers \DisplFile{clean(..)} method
\item Remove stored installer information object
\end{enumerate}

\subsection[sec:installation process fallback]{Fallback if package is not available}

This fallback works only for uninstall\supref{sec:uninstall process} process!

If the package is not available, this means the installer can not read wich installer class should be used.
It asked the user to use the basic operation, in this case, the BasicInstaller is used as installer to make the operations.

\section[sec:installer interfaces and classes]{Installer interfaces and classes}

\begin{figure}[h]
\begin{emp}[Installer interfaces and classes](20, 20)
Class.information("InstallationInformation")()(
	"+getVersion() : Version",
	"+getFiles() : array",
	"+getInstalledPackage() : Package",
	"+getTargetPackage() : Package");

Class.installer("Installer")()(
	"+install($objInfo : InstallationInformation, $objVersion : Version) throws InstallerException : void",
	"+upgrade($objInfo : InstallationInformation, $objTargetVersion : Version) throws InstallerException : void",
	"+downgrade($objInfo : InstallationInformation, $objTargetVersion : Version) throws InstallerException : void",
	"+reinstall($objInfo : InstallationInformation) throws InstallerException : void",
	"+uninstall($objInfo : InstallationInformation) throws InstallerException : void",	
	"+clean($objInfo : InstallationInformation) throws InstallerException : void");

Class.basicInstaller("BasicInstaller")()(
	"+removeFiles($objInfo : InstallationInformation) throws InstallerException : void",
	"+removeContaoFiles($objInfo : InstallationInformation) throws InstallerException : void",
	"+removeDataFiles($objInfo : InstallationInformation) throws InstallerException : void",
	"+removeUnusedFiles($objInfo : InstallationInformation) throws InstallerException : void",
	"+removeUnusedContaoFiles($objInfo : InstallationInformation) throws InstallerException : void",
	"+removeUnusedDataFiles($objInfo : InstallationInformation) throws InstallerException : void",
	"+copyFiles($objInfo : InstallationInformation) throws InstallerException : void",
	"+copyContaoFiles($objInfo : InstallationInformation) throws InstallerException : void",
	"+copyUserFiles($objInfo : InstallationInformation) throws InstallerException : void",
	"+removeSchema($objInfo : InstallationInformation) throws InstallerException : void",
	"+safePatchSchema($objInfo : InstallationInformation) throws InstallerException : void",
	"+fullPatchSchema($objInfo : InstallationInformation) throws InstallerException : void");

topToBottom(20)(installer, basicInstaller);
topToBottom(20)(information, installer);

drawObjects(information, installer, basicInstaller);

clink(dashedLink)(information, installer);
clink(transition)(basicInstaller, installer);
\end{emp}
\caption{Installer interfaces and classes}
\end{figure}

\subsection[sec:installerinformation]{InstallationInformation}

The InstallationInformation object, contains all informations about an installation.

\begin{itemize}
\item Installed Version
\item Installed Files
\item Corresponding Schemas
\item The installed package file (if available)
\item The target installation package file (if a new version should be installed or up-/downgraded)
\end{itemize}

The InstallationInformation object is modified from the Installer.
The Installer register all it changes of the extension to the InstallationInformation object, including version changes, file list changes and schema changes.

\section[sec:installer]{Installer}

The installer have to provide methods to install, upgrade, downgrade and uninstall an extension.

\begin{figure}[h]
\begin{emp}[Installer Interface](20, 20)
Class.information("InstallationInformation")()(
	"+getVersion() : Version",
	"+getFiles() : array",
	"+getInstalledPackage() : Package",
	"+getTargetPackage() : Package");

Class.installer("Installer")()(
	"+install($objInfo : InstallationInformation, $objVersion : Version) throws InstallerException : void",
	"+upgrade($objInfo : InstallationInformation, $objTargetVersion : Version) throws InstallerException : void",
	"+downgrade($objInfo : InstallationInformation, $objTargetVersion : Version) throws InstallerException : void",
	"+reinstall($objInfo : InstallationInformation, $objTargetVersion : Version) throws InstallerException : void",
	"+uninstall($objInfo : InstallationInformation) throws InstallerException : void",	
	"+clean($objInfo : InstallationInformation) throws InstallerException : void");

topToBottom(20)(information, installer);

drawObjects(information, installer);

clink(dashedLink)(information, installer);
\end{emp}
\caption{Installer interface}
\end{figure}

\DisplCode{install(\$objInfo, \$objVersion)} Installs the extension. It gets the target version to install as argument to make it easer for developers to create multi version installers.

\vspace{.5cm}
\DisplCode{upgrade(\$objInfo, \$objTargetVersion)} Upgrades the extension from one version to another version.

\vspace{.5cm}
\DisplCode{downgrade(\$objInfo, \$objTargetVersion)} Downgrades the extension from one version to another version.

\vspace{.5cm}
\DisplCode{reinstall(\$objInfo, \$objTargetVersion)} Reinstall the extension with a specific version. The target version does not have to be the same, as the installed version. This method can be used to downgrade to a lower version, that is not supported by the \DisplCode{downgrade(..)} method.

\vspace{.5cm}
\DisplCode{uninstall(\$objInfo)} Uninstalls the extension. By definition, this will only uninstall the CONTAO files. All user DATA and Schema Information and data will be kept.

\vspace{.5cm}
\DisplCode{clean(\$objInfo)} Completely clean the installation, removes all files and schema data.

\vspace{.5cm}
The installation process success if there is no exception thrown.

\newpage
\section[sec:basic installer]{Basic Installer}

The Basic Installer is the default implementation of the installer interface.
If the developer does not define an installer for the extension, this installer is used.
In most cases, if the developer needs to add custom install, upgrade, downgrade or uninstall code, the developer create a supimplementation of this class.

\begin{figure}[h]
\begin{emp}[Basic Installer](20, 20)
Class.basicInstaller("BasicInstaller")()(
	"+removeFiles($objInfo : InstallationInformation) throws InstallerException : void",
	"+removeContaoFiles($objInfo : InstallationInformation) throws InstallerException : void",
	"+removeDataFiles($objInfo : InstallationInformation) throws InstallerException : void",
	"+removeUnusedFiles($objInfo : InstallationInformation) throws InstallerException : void",
	"+removeUnusedContaoFiles($objInfo : InstallationInformation) throws InstallerException : void",
	"+removeUnusedDataFiles($objInfo : InstallationInformation) throws InstallerException : void",
	"+copyFiles($objInfo : InstallationInformation) throws InstallerException : void",
	"+copyContaoFiles($objInfo : InstallationInformation) throws InstallerException : void",
	"+copyUserFiles($objInfo : InstallationInformation) throws InstallerException : void",
	"+removeSchema($objInfo : InstallationInformation) throws InstallerException : void",
	"+safePatchSchema($objInfo : InstallationInformation) throws InstallerException : void",
	"+fullPatchSchema($objInfo : InstallationInformation) throws InstallerException : void");

topToBottom(20)(installer, basicInstaller);

drawObjects(installer, basicInstaller);

clink(transition)(basicInstaller, installer);
\end{emp}
\caption{Basic Installer Implementation}
\end{figure}

% remove methods
\vspace{.5cm}
\DisplCode{removeFiles(\$objInfo)} Removes all CONTAO and DATA files from the installation.

\vspace{.5cm}
\DisplCode{removeContaoFiles(\$objInfo)} Removes all CONTAO files from the installation.

\vspace{.5cm}
\DisplCode{removeDataFiles(\$objInfo)} Removes all DATA files from the installation.

\vspace{.5cm}
\DisplCode{removeUnusedFiles(\$objInfo)} Removes unused CONTAO and DATA files from the installation.

\vspace{.5cm}
\DisplCode{removeUnusedContaoFiles(\$objInfo)} Removes unused CONTAO files from the installation.

\vspace{.5cm}
\DisplCode{removeUnusedDataFiles(\$objInfo)} Removes unused DATA files from the installation.

% copy files methods
\vspace{.5cm}
\DisplCode{copyContaoFiles(\$objInfo)} Copy the files from CONTAO package directory into the installation.

\vspace{.5cm}
\DisplCode{copyDataFiles(\$objInfo)} Copy the files from DATA package directory into the installation.

% schema methods
\vspace{.5cm}
\DisplCode{removeSchema(\$objInfo)} Remove the complete schema and all containing data from the database.

\vspace{.5cm}
\DisplCode{safePatchSchema(\$objInfo)} Add missing or modify existing schema, but does not delete any existing schema information and data.

\vspace{.5cm}
\DisplCode{fullPatchSchema(\$objInfo)} Add missing or modify existing schema, and delete no more existing schema information and data.

\vspace{.5cm}
The basic implementation may look like this. Please note, this is an incomplete, pseudo implementation!

\begin{lstlisting}[caption=Basic Installer Implementation Code]
class BasicInstaller implements Installer
{
	public function install(InstallationInformation \$objInfo, Version \$objVersion)
	{
		$this->copyContaoFiles($objInfo);
		$this->copyDataFiles($objInfo);
		$this->safePatchSchema($objInfo);
	}
	
	public function upgrade(InstallationInformation $objInfo, Version $objToVersion)
	{
		$this->removeUnusedFiles($objInfo);
		$this->safePatchSchema($objInfo);
		$this->copyFiles($objInfo);
	}
	
	public function downgrade(InstallationInformation $objInfo, Version $objToVersion)
	{
		$objVersionCompare = $objInfo->getVersion()->compare($objToVersion);
		
		// check if the versions differ in more than the build number
		if (!$objVersionCompare->isAnotherBuild())
		{
			throw new DowngradeException("Downgrade to another major, minor or revision is not supported.");
		}
		
		$this->removeUnusedFiles($objInfo);
		$this->safePatchSchema($objInfo);
		$this->copyFiles($objInfo);
	}
	
	public function reinstall(InstallationInformation $objInfo, Version $objTargetVersion)
	{
		$this->removeContaoFiles($objInfo);
		$this->safePatchSchema($objInfo);
		$this->copyFiles($objInfo);
	}
	
	public function uninstall(InstallationInformation $objInfo)
	{
		$this->removeContaoFiles($objInfo);
	}
	
	public function clean(InstallationInformation $objInfo)
	{
		$this->removeFiles($objInfo);
		$this->removeSchema($objInfo);
	}
	
	public function removeFiles(InstallationInformation $objInfo)
	{
		$this->removeContaoFiles($objInfo);
		$this->removeDataFiles($objInfo);
	}
	
	public function removeContaoFiles(InstallationInformation $objInfo)
	{
		// remove files from the CONTAO package directory.
		// unregister removed files from $objInfo
	}
	
	public function removeDataFiles(InstallationInformation $objInfo)
	{
		// remove files from the DATA package directory.
		// unregister removed files from $objInfo
	}
	
	public function removeUnusedFiles(InstallationInformation $objInfo)
	{
		$this->removeUnusedContaoFiles($objInfo);
		$this->removeUnusedDataFiles($objInfo);
	}
	
	public function removeUnusedContaoFiles(InstallationInformation $objInfo)
	{
		// remove files from the CONTAO package directory that are unused.
		// unregister removed files from $objInfo
	}
	
	public function removeUnusedDataFiles(InstallationInformation $objInfo)
	{
		// remove files from the DATA package directory that are unused.
		// unregister removed files from $objInfo
	}
	
	public function copyFiles(InstallationInformation $objInfo)
	{
		$this->copyContaoFiles($objInfo);
		$this->copyUserFiles($objInfo);
	}
	
	public function copyContaoFiles(InstallationInformation $objInfo)
	{
		// copy files from the CONTAO package into the installation.
		// register installed files into $objInfo
	}
	
	public function copyUserFiles(InstallationInformation $objInfo)
	{
		// copy files from the DATA package into the installation.
		// register installed files into $objInfo
	}
	
	public function removeSchema(InstallationInformation $objInfo)
	{
		// remove the complete schema
	}
	
	public function safePatchSchema(InstallationInformation $objInfo)
	{
		// add missing schema information and update existing
	}
	
	public function fullPatchSchema(InstallationInformation $objInfo)
	{
		// add missing schema information, update existing and remove no more existing informations.
	}
}
\end{lstlisting}

\section[sec:custom installer code]{Custom Installer Code}

If a developer want's to add custom code, he can make a subclass of the Basic Installer.

\begin{lstlisting}[caption=Custom Installer Code Example]
class CustomInstaller extends BasicInstaller
{
	public function upgrade(InstallationInformation $objInfo, Version $objToVersion)
	{
		// add custom pre-upgrade code here
		
		parent::upgrade($objInfo, $objToVersion);
		
		// add custom post-upgrade code here
	}	
}
\end{lstlisting}

\subsection[sec:basic installer]{Basic Installer}

The Basic Installer is the default implementation of the installer interface.
If the developer does not define an installer for the extension, this installer is used.
In most cases, if the developer needs to add custom install, upgrade, downgrade or uninstall code, the developer create a supimplementation of this class.

\begin{figure}[h]
\begin{emp}[Basic Installer](20, 20)
\end{emp}
\caption{Basic Installer Implementation}
\end{figure}

% remove methods
\vspace{.5cm}
\DisplCode{removeFiles(\$objInfo)} Removes all CONTAO and DATA files from the installation.

\vspace{.5cm}
\DisplCode{removeContaoFiles(\$objInfo)} Removes all CONTAO files from the installation.

\vspace{.5cm}
\DisplCode{removeDataFiles(\$objInfo)} Removes all DATA files from the installation.

\vspace{.5cm}
\DisplCode{removeUnusedFiles(\$objInfo)} Removes unused CONTAO and DATA files from the installation.

\vspace{.5cm}
\DisplCode{removeUnusedContaoFiles(\$objInfo)} Removes unused CONTAO files from the installation.

\vspace{.5cm}
\DisplCode{removeUnusedDataFiles(\$objInfo)} Removes unused DATA files from the installation.

% copy files methods
\vspace{.5cm}
\DisplCode{copyContaoFiles(\$objInfo)} Copy the files from CONTAO package directory into the installation.

\vspace{.5cm}
\DisplCode{copyDataFiles(\$objInfo)} Copy the files from DATA package directory into the installation.

% schema methods
\vspace{.5cm}
\DisplCode{removeSchema(\$objInfo)} Remove the complete schema and all containing data from the database.

\vspace{.5cm}
\DisplCode{safePatchSchema(\$objInfo)} Add missing or modify existing schema, but does not delete any existing schema information and data.

\vspace{.5cm}
\DisplCode{fullPatchSchema(\$objInfo)} Add missing or modify existing schema, and delete no more existing schema information and data.

\vspace{.5cm}
The basic implementation may look like this. Please note, this is an incomplete, pseudo implementation!

\begin{lstlisting}[caption=Basic Installer Implementation Code]
class BasicInstaller implements Installer
{
	public function install(InstallationInformation \$objInfo, Version \$objVersion)
	{
		$this->copyContaoFiles($objInfo);
		$this->copyDataFiles($objInfo);
		$this->safePatchSchema($objInfo);
	}
	
	public function upgrade(InstallationInformation $objInfo, Version $objToVersion)
	{
		$this->removeUnusedFiles($objInfo);
		$this->safePatchSchema($objInfo);
		$this->copyFiles($objInfo);
	}
	
	public function downgrade(InstallationInformation $objInfo, Version $objToVersion)
	{
		$objVersionCompare = $objInfo->getVersion()->compare($objToVersion);
		
		// check if the versions differ in more than the build number
		if (!$objVersionCompare->isAnotherBuild())
		{
			throw new DowngradeException("Downgrade to another major, minor or revision is not supported.");
		}
		
		$this->removeUnusedFiles($objInfo);
		$this->safePatchSchema($objInfo);
		$this->copyFiles($objInfo);
	}
	
	public function reinstall(InstallationInformation $objInfo, Version $objTargetVersion)
	{
		$this->removeContaoFiles($objInfo);
		$this->safePatchSchema($objInfo);
		$this->copyFiles($objInfo);
	}
	
	public function uninstall(InstallationInformation $objInfo)
	{
		$this->removeContaoFiles($objInfo);
	}
	
	public function clean(InstallationInformation $objInfo)
	{
		$this->removeFiles($objInfo);
		$this->removeSchema($objInfo);
	}
	
	public function removeFiles(InstallationInformation $objInfo)
	{
		$this->removeContaoFiles($objInfo);
		$this->removeDataFiles($objInfo);
	}
	
	public function removeContaoFiles(InstallationInformation $objInfo)
	{
		// remove files from the CONTAO package directory.
		// unregister removed files from $objInfo
	}
	
	public function removeDataFiles(InstallationInformation $objInfo)
	{
		// remove files from the DATA package directory.
		// unregister removed files from $objInfo
	}
	
	public function removeUnusedFiles(InstallationInformation $objInfo)
	{
		$this->removeUnusedContaoFiles($objInfo);
		$this->removeUnusedDataFiles($objInfo);
	}
	
	public function removeUnusedContaoFiles(InstallationInformation $objInfo)
	{
		// remove files from the CONTAO package directory that are unused.
		// unregister removed files from $objInfo
	}
	
	public function removeUnusedDataFiles(InstallationInformation $objInfo)
	{
		// remove files from the DATA package directory that are unused.
		// unregister removed files from $objInfo
	}
	
	public function copyFiles(InstallationInformation $objInfo)
	{
		$this->copyContaoFiles($objInfo);
		$this->copyUserFiles($objInfo);
	}
	
	public function copyContaoFiles(InstallationInformation $objInfo)
	{
		// copy files from the CONTAO package into the installation.
		// register installed files into $objInfo
	}
	
	public function copyUserFiles(InstallationInformation $objInfo)
	{
		// copy files from the DATA package into the installation.
		// register installed files into $objInfo
	}
	
	public function removeSchema(InstallationInformation $objInfo)
	{
		// remove the complete schema
	}
	
	public function safePatchSchema(InstallationInformation $objInfo)
	{
		// add missing schema information and update existing
	}
	
	public function fullPatchSchema(InstallationInformation $objInfo)
	{
		// add missing schema information, update existing and remove no more existing informations.
	}
}
\end{lstlisting}

\section[sec:custom installer code]{Custom Installer}

If a developer want's to add custom code, he can make a subclass of the BasicInstaller.

\begin{lstlisting}[caption=Custom Installer Code Example]
class CustomInstaller extends BasicInstaller
{
	public function upgrade(InstallationInformation $objInfo, Version $objToVersion)
	{
		// add custom pre-upgrade code here
		
		parent::upgrade($objInfo, $objToVersion);
		
		// add custom post-upgrade code here
	}	
}
\end{lstlisting}

Or the developer write his complete own installer. In this case the developer have to implement the Installer interface.


% UML support
\end{empfile}

% Seitennummerierung -----------------------------------------------------------
%   Nach dem Hauptteil werden die Seiten in großen römischen Ziffern 
%   nummeriert.
% ------------------------------------------------------------------------------
\pagenumbering{Roman}

% Anhang -----------------------------------------------------------------------
%   Die Inhalte des Anhangs werden analog zu den Kapiteln inkludiert.
%   Dies geschieht in der Datei "Appendix.tex".
% ------------------------------------------------------------------------------
\begin{appendix}
    \clearpage
    \chapter{Appendix}
    \label{sec:Appendix}
    \pagenumbering{roman}
    % Rand der Aufzählungen in Tabellen anpassen
    \setdefaultleftmargin{0em}{}{}{}{}{}
    \section[sec:package information file schema]{PACKAGE information file schema}
\xmllisting{lst:extension.xsd}{PACKAGE information file schema}{../extension.xsd}

\section[sec:package information file example]{PACKAGE information file example}
\xmllisting{lst:extension.xml}{PACKAGE information file example}{../extension.xml}

\end{appendix}

% Abkürzungsverzeichnis --------------------------------------------------------
\nomenclature{API}{Application Programming Interface}
\nomenclature{ARIS}{Architektur integrierter Informationssysteme}
\nomenclature{BPR}{Business Process Reengineering}
\nomenclature{eEPK}{erweiterte Ereignisgesteuerte Prozesskette}
\nomenclature{EPK}{Ereignisgesteuerte Prozesskette}
\nomenclature{JMS}{Java Message Service}
\nomenclature{SDK}{Software Development Kit}
\nomenclature{URI}{Uniform Resource Identifier}
\nomenclature{URL}{Uniform Resource Locator}
\nomenclature{URN}{Uniform Resource Name}
\nomenclature{W3C}{World Wide Web Consortium}
\nomenclature{XML}{Extensible Markup Language}
\nomenclature{XPath}{XML Path Language}
\nomenclature{XSL}{Extensible Stylesheet Language}
\nomenclature{XSLT}{XSL Transformations}

\glossary{name={HOOK}, description={Ein Einsprungpunkt bei dem eine Erweiterung in den Ablauf einer Funktion eingehängt wird.}}
% für korrekte Überschrift in der Kopfzeile
\clearpage\markboth{\nomname}{\nomname} 
\printnomenclature
\label{sec:Glossar}

\listoffigures % Abbildungsverzeichnis
\listoftables % Tabellenverzeichnis
\renewcommand{\lstlistlistingname}{List of Listings}
\lstlistoflistings % Listings-Verzeichnis

% Literaturverzeichnis ---------------------------------------------------------
%   Das Literaturverzeichnis wird aus der BibTeX-Datenbank "Bibliographie.bib"
%   erstellt.
% ------------------------------------------------------------------------------
\bibliography{Bibliographie} % Aufruf: bibtex RFCCC1
\bibliographystyle{natdin} % DIN-Stil des Literaturverzeichnisses

% Index ------------------------------------------------------------------------
%   Zum Erstellen eines Index, die folgende Zeile auskommentieren.
% ------------------------------------------------------------------------------
\printindex
\printglossaries

\end{document}
